\section{The \method Baseline}
\label{sec:method}

\begin{figure*}[t]
\centering
\resizebox{0.98\textwidth}{!}{%
\begin{tikzpicture}[
  font=\scriptsize,
  node distance=12mm and 6mm,
  box/.style={draw,rounded corners=2pt,minimum height=10mm,text width=31mm,align=center,inner xsep=4pt,inner ysep=3pt},
  data/.style={box,fill=lightblue,draw=eventblue},
  audit/.style={box,fill=lightorange,draw=eventorange},
  model/.style={box,fill=lightgreen,draw=eventgreen},
  loss/.style={box,fill=red!5,draw=eventred},
  arrow/.style={-Latex,thick,draw=black!70}
]
\node[data] (pair) {unique context $x_i^g$ + crops $\{x_{ij}^l\}$\\content / flight / site group};
\node[audit,right=of pair] (card) {sample graph + evidence cards\\events, factors, assignments};
\node[audit,right=of card] (caption) {label-agnostic description\\grounded rewrite + audit};
\node[data,right=of caption] (tiers) {versioned benchmark\\Core / Tail / Catalog};

\node[model,below=of pair] (student) {Qwen3-VL-8B\\frozen vision + LoRA};
\node[model,right=of student] (views) {context--crop edge scores\\view / evidence dropout};
\node[model,right=of views] (score) {permutation-invariant set score\\event set + crop assignment};
\node[data,right=of score] (eval) {group-disjoint evaluation\\hard negatives + calibration};

\node[loss,text width=70mm,below=7mm of $(views.south)!0.5!(score.south)$,anchor=north] (losses)
{$\mathcal{L}_{\mathrm{gen}}$ + taxonomy-neighbor margin + counterfactual unlikelihood};

\draw[arrow] (pair) -- (card);
\draw[arrow] (card) -- (caption);
\draw[arrow] (caption) -- (tiers);
\draw[arrow] (pair.south) -- (student.north);
\draw[arrow] (student) -- (views);
\draw[arrow] (views) -- (score);
\draw[arrow] (score) -- (eval);
\draw[arrow] (card.south) -- (views.north);
\draw[arrow] (tiers.south) -- (eval.north);
\draw[arrow] ([xshift=-13mm]losses.north) -- ([xshift=5mm]views.south);
\draw[arrow] ([xshift=13mm]losses.north) -- ([xshift=-5mm]score.south);

\node[font=\footnotesize\bfseries,text=eventorange,above=2mm of pair.north west,anchor=south west] {Auditable data engine};
\node[font=\footnotesize\bfseries,text=eventgreen,above=2mm of student.north west,anchor=south west] {\method baseline};
\end{tikzpicture}}
\caption{\dataset separates an auditable data engine (top) from \method/CLEAR-Set and frozen evaluation (bottom). Exact-context grouping converts pair rows into variable-size sample graphs; edge scores are aggregated without crop order while retaining evidence assignments.}
\label{fig:method}
\end{figure*}


\subsection{Dual-View Structured Prediction}

\method instantiates a deliberately reproducible baseline rather than a new foundation model. A Qwen3-VL-8B-Instruct student~\cite{bai2025qwen3vl} receives the global context $x^g$ and local evidence $x^l$ as two ordered images with explicit \texttt{CONTEXT} and \texttt{EVIDENCE} markers. The pair-level target is
\begin{equation}
\begin{aligned}
s=[&\texttt{events}:Y;\ \texttt{factors}:\{(d,e,s,r,c)\}_{y\in Y};\\
   &\texttt{evidence}:q;\ \texttt{uncertain}:u],
\end{aligned}
\end{equation}
where $q$ is a short visible-evidence statement and $u$ is a binary abstention field. Language-model and projector LoRA modules are trained while the vision tower and base language weights remain frozen. Single-view baselines use the same prompt and token budget.

Training uses verified event phrases, factors, and audited evidence statements. The model is never given a candidate list in the free-generation setting. In the closed-set setting, the normalized conditional log-likelihood of each canonical event definition gives a reproducible label-wise score
\begin{equation}
 S_{\theta}(y\mid x^g,x^l)=\frac{1}{|a_y|}\sum_{t\in a_y}
 \log p_{\theta}(a_{y,t}\mid a_{y,<t},x^g,x^l),
\end{equation}
where $a_y$ is the canonical label plus its compact visual definition. Temperature-scaled label scores are thresholded independently, allowing legitimate co-occurrence rather than forcing a softmax winner.

\subsection{Permutation-Invariant Set Extension}

For a context set $G_i$ with $m_i$ crops, \method evaluates every context--crop edge and aggregates evidence without crop order:
\begin{equation}
 S_{\theta}^{\mathrm{set}}(y\mid G_i)=
 \tau\log\sum_{j=1}^{m_i}\exp\!\left(S_{\theta}(y\mid x_i^g,x_{ij}^l)/\tau\right).
\end{equation}
Binary cross-entropy over $Y_i$ trains set prediction, while the maximizing edge supplies a transparent evidence assignment $\hat{A}_i$. We compare log-sum-exp with max pooling and a parameter-matched attention pool. This \emph{CLEAR-Set} extension adds no crop-order encoding and is scored only on adjudicated multi-evidence groups.

\subsection{Graph-Structured Supervision}

Token-level cross-entropy alone rewards fluent background description and ignores which confusion was resolved. We retain label-span weighting from conventional grounded caption tuning, but add two graph-aware terms. For each positive $y\in Y$, a neighbor $y^{-}\in\mathcal{N}(y)\setminus Y$ shares the entity or scene while differing in one decision-critical factor. A margin loss requires
\begin{equation}
 \mathcal{L}_{\mathrm{nbr}}=
 \max\!\left(0,m-S_{\theta}(y\mid x^g,x^l)+S_{\theta}(y^{-}\mid x^g,x^l)\right).
\end{equation}
We also construct a counterfactual evidence statement $q^{-}$ by changing only the disputed factor. Its token probability is penalized with an unlikelihood objective
\begin{equation}
 \mathcal{L}_{\mathrm{cf}}=-\sum_{t\in q^{-}}\log\left(1-p_{\theta}(q^{-}_t\mid q^{-}_{<t},x^g,x^l)\right).
\end{equation}
The full objective is
\begin{equation}
 \mathcal{L}=\mathcal{L}_{\mathrm{gen}}+
 \lambda_{\mathrm{nbr}}\mathcal{L}_{\mathrm{nbr}}+
 \lambda_{\mathrm{cf}}\mathcal{L}_{\mathrm{cf}},
\end{equation}
where $\mathcal{L}_{\mathrm{gen}}$ upweights canonical label tokens by $\gamma$. Ablations remove each term and replace graph neighbors with random negatives to test whether gains come from decision-boundary structure.

\subsection{Context--Evidence Reliance}

Joint input can still ignore one view. We therefore train with controlled view dropout: context-only, evidence-only, and joint examples appear with fixed probabilities. The target stays unchanged only when the remaining view satisfies the event's required evidence factors; otherwise \texttt{uncertain} is set true. At evaluation, four conditions isolate the perception pipeline: context-only, oracle crop, context plus oracle crop, and context plus a detector/retriever proposal. The oracle-to-proposal gap measures localization error, while the joint-to-best-single-view gap measures genuine complementarity.

Accepting two images does not prove cross-view synthesis~\cite{guo2026cvbench}. We therefore add three necessity diagnostics on adjudicated graphs. \emph{Single-view insufficiency success} counts items that are correct jointly but incorrect from either view alone. \emph{Assigned-evidence deletion} removes the human-linked crop for an event and compares the score drop with deletion of a size-matched random crop. Finally, a crop-budget curve reports set F1 as the model admits its top-$k$ evidence edges. A method that cherry-picks one easy view will have a small assigned-versus-random deletion gap and a flat joint advantage even if its headline set F1 is high.

\subsection{Confidence and Selective Prediction}

Closed-set label scores are temperature-scaled on validation groups. The model abstains when no label exceeds its validation-selected threshold or when factor outputs violate the \eventgraph. We report classwise expected calibration error and risk--coverage curves; abstention is never counted as correct. Free-generation outputs are mapped to the taxonomy by exact canonical/synonym matching, followed by a frozen embedding tie-breaker, with all unmatched phrases counted as errors.
